\subsection{Strict maps}

PROPOSITION 8. --- Let

\[
\begin{array}{c} \mathrm {X^ {\prime}} \xrightarrow {f ^ {\prime}} \mathrm{X} \\ \Big \downarrow_ {p ^ {\prime}} \\ \mathrm {B^ {\prime}} \xrightarrow {f} \mathrm{B} \end{array}
\]

a Cartesian square. If the map p is open (resp. is proper, resp. has a continuous section in a neighborhood of every point), then so is \(p'\).

By Remark 2, I, p. 16, it suffices to prove the proposition for Cartesian squares of the following type:

\[
\begin{array}{c c} \overline {{p}} ^ {1} (\mathrm {A}) \xrightarrow {j} \mathrm{X} & \mathrm {A} \times \mathrm{F} \xrightarrow {i} \mathrm{X} \times \mathrm{F} \\ \Big \downarrow_ {p _ {\mathrm {A}}} & \Big \downarrow_ {p \times \mathrm {Id}_{\mathrm {F}}} \\ \mathrm {A} & \xrightarrow {i} \mathrm {B} \end{array} \quad\text {and} \quad \begin{array}{c} \mathrm {X} \times \mathrm {F} \xrightarrow {j} \mathrm {X} \\ \Big \downarrow_ {p \times \mathrm {Id}_{\mathrm {F}}} \\ \mathrm {B} \times \mathrm {F} \xrightarrow {i} \mathrm {B} \end{array}
\]

where F is a topological space, A is a subspace of B, and i and j are the canonical injections. If the map p is open, the maps \(p_{A}\) and \(p \times Id_{F}\) are open (TG, I, p. 30, prop. 2 and TG, I, p. 34, prop. 8). If the map p is proper, then the maps \(p_{A}\) and \(p \times Id_{F}\) are proper (TG, I, p. 72, prop. 3 and TG, I, p. 72, def. 1). If U is an open subset of B and \(s\colon\mathrm{U}\to\mathrm{X}\) is a continuous section of p on U, the map \(s\mid(\mathrm{A}\cap\mathrm{U})\) is a continuous section of \(p_{\mathrm{A}}\) onto \(\mathrm{A}\cap\mathrm{U}\), and the map \(s\times\mathrm{Id}_{\mathrm{F}}\) is a continuous section of \(p\times\mathrm{Id}_{\mathrm{F}}\) onto \(\mathrm{U}\times\mathrm{F}\).

Remark 1. --- With the notation of Proposition 8, if \(p\) is a closed map, it is not necessarily the case that \(p'\) is closed (cf. TG, I, p. 32, prop. 3). However, if \(p\) is a map with a continuous section, \(p'\) also has a continuous section.

DEFINITION 5. --- Let X and Y be topological spaces and let f: X → Y be a map. Let R be the equivalence relation associated with f and

\[
\mathrm{X} \rightarrow \mathrm{X} / \mathrm{R} \xrightarrow {g} f (\mathrm{X}) \rightarrow \mathrm{Y}
\]

the canonical decomposition of \(f\) (E, II, p. 44). We say that the map \(f\) is strict if \(g\) is a homeomorphism, when \(X/R\) is equipped with the quotient topology and \(f(X)\) is equipped with the topology induced by that of \(Y\).

A strict map is continuous.

Recall (TG, I, p. 22, prop. 8) that for a map \(f\) to be strict, it is necessary and sufficient that \(f\) be continuous and that, for every subset \(A\) of \(X\) which is open (resp. closed) and saturated, \(f(A)\) be open (resp. closed) in \(f(X)\).

Examples. --- 1) The composition of two strict maps is not necessarily strict. Indeed, every continuous map \(f \colon X \to Y\) is the composite of the map \(\mathrm{pr}_2 \colon X \times Y \to Y\) and the map \(x \mapsto (x, f(x))\) from \(X\) into \(X \times Y\), both of which are strict (TG, I, p. 26, prop. 5 and TG, I, p. 25, cor. 2). On the other hand, the composite of two strict and injective (resp. surjective) maps is a strict map.

\begin{enumerate}
\def\labelenumi{\arabic{enumi})}
\setcounter{enumi}{1}
\item
  A continuous map that is open, or closed, or has a continuous section, is strict. This follows from Prop. 3 of TG, I, p. 32 and Prop. 9 of TG, I, p. 22.
\item
  For a continuous homomorphism from one topological group to another to be a strict morphism (TG, III, p. 16, def. 1), it is necessary and sufficient that it be a strict map in the sense of Definition 5.
\end{enumerate}

PROPOSITION 9. --- Let X, Y, and Z be topological spaces. Let \(f\colon X \to Y\) be a continuous surjective map and \(g\colon Y \to Z\) a map.

\begin{enumerate}
\def\labelenumi{\alph{enumi})}
\item
If f is strict and if g∘f is continuous, then the map g is continuous.
\item
If g is continuous and if g∘f is strict, then g is strict.
\item
If f and g \(\circ\) f are strict, then g is strict.
\end{enumerate}

Let us prove assertion a). Let R denote the relation associated with f in X; by hypothesis, the map from X/R onto Y induced by f by passing to the quotient is a homeomorphism. The first assertion then follows from Proposition 6 of I, p. 21.

Let us prove b). Let B be a closed saturated subset of Y for the relation defined by g, and let \(A = f^{-1}(B)\). Since f is continuous, A is closed in X, and A is saturated for the equivalence relation defined by \(g \circ f\). Since \(g \circ f\) is strict and f is surjective, \(g(B) = g \circ f(A)\) is then closed in Z. The continuous map g is therefore strict.

Assertion c) follows immediately from assertions a) and b).

PROPOSITION 10. --- Let X be a topological space and let R be an equivalence relation on X. Let Y be a locally compact topological space. Let S be the equivalence relation on X × Y which is the product of the equivalence relation R on X and the equality relation on Y. The canonical bijection (X × Y)/S → (X/R) × Y is a homeomorphism.

Recall that if U and V are topological spaces, \(\mathcal{C}_{c}(\mathrm{U};\mathrm{V})\) denotes the set of continuous maps from U to V, endowed with the compact-open topology (TG, X, p. 26, def. 1).

Let \(p \colon \mathrm{X} \to \mathrm{X} / \mathrm{R}\) and \(q \colon \mathrm{X} \times \mathrm{Y} \to (\mathrm{X} \times \mathrm{Y}) / \mathrm{S}\) be the canonical surjections. Let \(g \colon (\mathrm{X} \times \mathrm{Y}) / \mathrm{S} \to (\mathrm{X} / \mathrm{R}) \times \mathrm{Y}\) be the canonical bijection. It is continuous; let us denote its inverse by \(h\) and prove that it is continuous.

The map \(i\colon X \to \mathcal{C}_{c}(Y;X \times Y)\) such that, for every \(x \in X\), \(i(x)\) is the map defined by \(y \mapsto (x,y)\), is continuous (TG, X, p. 28, Theorem 3). The map \(\widetilde{q}\colon \mathcal{C}_{c}(Y;X \times Y) \to \mathcal{C}_{c}(Y;(X \times Y)/S)\), which associates to a continuous map \(\varphi\) the map \(q \circ \varphi\), is continuous (TG, X, p. 29, Proposition 9). Consequently, the map \(\widetilde{q} \circ i\colon X \to \mathcal{C}_{c}(Y;(X \times Y)/S)\) is continuous. It is compatible with the equivalence relation R; the unique map \(j\colon X/R \to \mathcal{C}_{c}(Y;(X \times Y)/S)\) such that \(j \circ p = \widetilde{q} \circ i\) is therefore continuous. We have \(h(\xi,y) = j(\xi)(y)\) for every pair \((\xi,y) \in (X/R) \times Y\). Since Y is locally compact, it follows from TG, X, p. 28, Theorem 3 that \(h\) is continuous.

The conclusion of Proposition 10 is no longer necessarily valid if Y is not locally compact (TG, I, p. 96, exerc. 6).

COROLLARY. --- Let X and Y be topological spaces and let \(f: X \to Y\) be a continuous map. Let T be a locally compact topological space. If the map f is strict, then so is the map \(f \times Id_{T}\) from \(X \times T\) into \(Y \times T\).

